\chapter*{RINGKASAN}
\addcontentsline{toc}{chapter}{RINGKASAN}

Interoperabilitas data kesehatan memerlukan penyelarasan makna istilah klinis dengan terminologi standar. Variasi istilah lokal, singkatan, dan campuran bahasa dalam narasi Indonesia menjadi tantangan pemetaan. CDE-Mapper telah mendukung berbagai kosakata, tetapi penerapannya pada narasi Indonesia memerlukan adaptasi. Penelitian ini mengembangkan dan mengevaluasi adaptasi CDE-Mapper untuk pemetaan terpadu ke SNOMED CT, LOINC, dan ICD-10; ekstraksi banyak \textit{clinical data elements} (CDE) dari teks panjang; serta adaptasi \textit{retrieval ensemble} terhadap bahasa Indonesia. Ketiga terminologi dipilih sesuai domain dan konteks dan kebutuhan representasi informasi.

Metode mengintegrasikan normalisasi, segmentasi, ekstraksi multi-entitas berbasis \textit{large language model} (LLM), pemilihan terminologi target, serta dekomposisi CDE dengan mempertahankan atribut, negasi, dan temporalitas. Pencarian menggunakan \textit{hybrid retrieval} yang menggabungkan jalur leksikal dan semantik, dilanjutkan penyaringan dan pemeringkatan ulang berdasarkan domain, atribut, granularitas, serta konsistensi penilaian LLM. Kasus di bawah ambang penerimaan menghasilkan keputusan \textit{abstain} untuk ditinjau ahli. Pemetaan tervalidasi disimpan dalam \textit{knowledge reservoir} dengan memperhatikan konteks dan versi terminologi.

Data direncanakan berasal dari rekam medis elektronik lokal dan konsultasi kesehatan daring yang dianonimkan. Tenaga medis menyusun \textit{gold standard}, dengan pemeriksaan kesepakatan antar-anotator dan adjudikasi. Data pengembangan dipisahkan dari data uji; jawaban acuan data uji tidak digunakan untuk mengisi reservoir. Evaluasi memisahkan kualitas ekstraksi, pemetaan dengan masukan entitas yang sama, dan kinerja \textit{end-to-end}. Pembanding mencakup pencocokan leksikal, LLM tanpa RAG, dan CDE-Mapper tanpa adaptasi bahasa Indonesia. Pengujian meliputi ablasi komponen dan ketahanan terhadap variasi bahasa, menggunakan akurasi pemetaan, kualitas peringkat kandidat, presisi, \textit{recall}, F1, cakupan, konsistensi, waktu inferensi, serta analisis pertukaran akurasi dan cakupan keputusan otomatis.

Luaran yang diharapkan berupa kerangka adaptasi dan prototipe pemetaan yang dapat ditelusuri untuk mendukung interoperabilitas semantik. Peningkatan akurasi dan efisiensi akan ditentukan melalui evaluasi.

\textbf{Kata kunci}: CDE-Mapper, teks klinis bahasa Indonesia, pemetaan multi-terminologi, LLM, \textit{hybrid retrieval}.